\documentclass[11pt]{article}
\usepackage[T1]{fontenc}
\usepackage{lmodern}
\usepackage[margin=1.05in]{geometry}
\usepackage{amsmath,amssymb,amsthm}
\usepackage[hidelinks]{hyperref}

% Optional author details are separate from the project contact below.
\newcommand*{\AuthorName}{}
\newcommand*{\AuthorEmail}{}
\newcommand{\byline}{\AuthorName\ifx\AuthorEmail\empty\else
  \ifx\AuthorName\empty\else\\\fi
  {\small\href{mailto:\AuthorEmail}{\expandafter\nolinkurl\expandafter{\AuthorEmail}}}\fi}
\newtheorem{theorem}{Theorem}[section]
\newtheorem{lemma}[theorem]{Lemma}
\newtheorem{corollary}[theorem]{Corollary}
\title{Large prime divisors of Fermat congruences\\and semiprime pseudoprimes}
\author{\small \href{mailto:admin@cosmosgeek.org}{admin@cosmosgeek.org}\qquad \href{mailto:b321014@iiit-bh.ac.in}{b321014@iiit-bh.ac.in}}
\date{}

\begin{document}
\maketitle
\begin{abstract}
We count the odd integers \(n\le X\) with exactly two prime factors,
counted with multiplicity, that satisfy \(2^{n-1}\equiv1\pmod n\).
Their number is \(O(X^{3/4}/(\log X)^{5/4})\). For distinct
factors \(r,s\), we group the pairs by \(\gcd(r-1,s-1)\).
A bound on the number of prime divisors of a single Mersenne number
handles small gcds, and two applications of the Brun--Titchmarsh
inequality handle large gcds. We also bound the number of composite
\(n\le X\) having a prime \(q>Y\) that divides both \(n\) and
\(2^{n-1}-1\) by \(O((X/Y)^2/\log Y)\). A large prime factor
of \(n-1\) gives a further bound for prime powers common to
\(n\) and \(a^{n-1}-1\) at any fixed base \(a\ge2\). These estimates quantify
the composite contribution to two possible counts of twin primes.
We also estimate false positives in simultaneous multiplicative-order tests and separate three explicit background terms in a shifted logarithmic gcd moment. For two independent random bases, a rank-zero elliptic curve bounds intersections of composite templates, giving a squared passing weight of $O_\epsilon(X^{3/7+\epsilon})$ on the specified rough composite pairs. An upper sieve for the mixed prime--semiprime rows gives an untrimmed squared-weight error of $O_\epsilon(\sqrt X/(\log X)^{3-\epsilon})$. A lower bound large enough to force infinitely many twin primes remains open.
\end{abstract}

\section{Semiprime pseudoprimes}

An odd composite integer \(n\) is a base-\(2\) Fermat pseudoprime if
\(n\mid 2^{n-1}-1\). Write \(Q_2(X)\) for the number of such integers
\(n\le X\) that are products of two primes, allowing equal factors.
All logarithms are natural except where the base is indicated. The
constants in \(O(\cdot)\) and \(\ll\) are absolute unless a
dependence is indicated.

\begin{theorem}\label{thm:semiprime}
For \(X\ge3\),
\[
             Q_2(X)\ll \frac{X^{3/4}}{(\log X)^{5/4}}.
\]
\end{theorem}

Luca, Ralaivaosaona and Jim\'enez Urroz proved
\(Q_2(X)\ll X^{4/5}/(\log X)^{2/5}\) in
\cite[Theorem~3.1]{LRJ}. Their result concerns products up to \(X\)
with unrestricted factor sizes. Their separate estimate for factors
in a common dyadic interval uses a different counting parameter.
Theorem~\ref{thm:semiprime} sharpens their bound for the unrestricted
count. We use their parametrization by the common divisor of the two
shifted prime factors and count the small values of that parameter
one at a time.

We use \(\varphi\) for Euler's totient function, \(\mu\) for the
M\"obius function and \(w(h)=h/\varphi(h)\).
The following two estimates keep track of both primality conditions
in the large-gcd range.

\begin{lemma}\label{lem:BT}
For all positive integers \(h\) and real \(T\ge2\),
\[
 \#\{1\le t\le T: ht+1\text{ is prime}\}
       \ll w(h)\frac{T}{\log(2T)}.
\]
For positive integers \(a\) and real \(M\ge2\),
\[
 \sum_{\substack{M<g\le2M\\ ag+1\ \mathrm{prime}}}w(g)
        \ll w(a)\frac{M}{\log(2M)}.
\]
In addition \(\sum_{a\le Z} w(a)/a\ll\log(2Z)\) for \(Z\ge1\).
\end{lemma}

\begin{proof}
The first inequality is the Brun--Titchmarsh inequality applied
to the class \(1\bmod h\) up to \(hT+1\) \cite[Theorem~2, (1.10)]{MV}.
For the second, use
\[
            w(g)=\sum_{d\mid g}\frac{\mu^2(d)}{\varphi(d)}.
\]
If \(d\le\sqrt M\), write \(g=dt\) and apply the first estimate
with \(h=ad\) and \(T=2M/d\). Since \(w(ad)\le w(a)w(d)\),
the total contribution of these \(d\) is at most a constant times
\[
 \frac{M w(a)}{\log(2M)}
      \sum_{d\ge1}\frac{\mu^2(d)}{\varphi(d)^2}.
\]
The series converges by its Euler product. For \(d>\sqrt M\),
count all multiples of \(d\) up to \(2M\). The inequality
\(\varphi(d)\ge\sqrt{d/2}\), obtained prime power by prime power,
bounds their contribution by
\(2M\sum_{d>\sqrt M}1/(d\varphi(d))\ll M^{3/4}\).
This gives the second assertion. Finally, expanding \(w(a)\) in the
same way and summing the harmonic series gives
\[
  \sum_{a\le Z}\frac{w(a)}{a}
       \le (1+\log Z)\sum_{d\ge1}\frac{\mu^2(d)}{d\varphi(d)}
       \ll\log(2Z).
\]
\end{proof}

\begin{proof}[Proof of Theorem~\ref{thm:semiprime}]
Prime squares contribute at most \(\sqrt X\). For the remaining
integers write \(n=rs\), with \(r<s\) odd primes, and set
\(g=\gcd(r-1,s-1)\). The order of \(2\) modulo \(r\) divides
both \(r-1\) and \(rs-1\), whose gcd is \(g\). Exchanging \(r\)
and \(s\) gives
\[
                  r,s\mid 2^g-1,\qquad r,s\equiv1\pmod g.
\]
For fixed \(g\ge2\), let \(t_g\) count the prime divisors of
\(2^g-1\) congruent to \(1\bmod g\). Their product is at least
\((g+1)^{t_g}\) and at most \(2^g-1\), whence
\(t_g\ll g/\log(2g)\). The pairs with \(g\le G\) therefore
contribute
\begin{equation}\label{eq:small}
     \sum_{2\le g\le G}t_g^2\ll \frac{G^3}{(\log(2G))^2}.
\end{equation}
For example, splitting the sum at \(\sqrt G\) gives the bound directly.

Take \(G=(X\log X)^{1/4}\) and assume \(X\) is large enough
that \(G<X^{3/8}\). For \(G<g\le X^{3/8}\), divide the range into
intervals \(M<g\le2M\), starting at \(M=G\) and truncating the
last one. Put \(r=ag+1\), \(s=bg+1\). Then
\[
       a\le \sqrt X/g,\qquad b\le B:=X/(ag^2),\qquad
       B\ge \sqrt X/g\ge X^{1/8}.
\]
For each eligible \(g,a\), Lemma~\ref{lem:BT} bounds the choices
of \(b\) for which \(bg+1\) is prime by
\(\ll Xw(g)/(ag^2\log X)\). Summing first over \(g\), while
retaining the requirement that \(ag+1\) be prime, bounds the number
of pairs in this block by
\begin{align*}
   &\ll \frac{X}{M^2\log X}
       \sum_{a\le\sqrt X/M}\frac1a
       \sum_{\substack{M<g\le2M\\ag+1\ \mathrm{prime}}}w(g)\\
   &\ll \frac{X}{M\log X\log(2M)}
               \sum_{a\le\sqrt X/M}\frac{w(a)}a
     \ll \frac{X}{M\log X}.
\end{align*}
Here \(\log(2M)\gg\log X\) because \(M\ge G\). The sum over
these blocks is \(O(X/(G\log X))\).

For \(g>X^{3/8}\), count every pair of positive integers
\(a,b\) with \(ab\le X/g^2\). There are at most
\((X/g^2)(1+\log X)\) such pairs. Summing over \(g\) gives
\(O(X^{5/8}\log X)\). At our choice of \(G\), this last
term and the prime squares are absorbed by the claimed bound,
and both \eqref{eq:small} and \(X/(G\log X)\) have the
required order.
\end{proof}

\section{One large common prime divisor}

A single prime factor of \(n\) can satisfy the Fermat congruence
even when \(n\) fails it. The following estimate applies to all
composites, with any number of prime factors.

\begin{theorem}\label{thm:common}
Let \(2\le Y\le X\). The number of pairs \((n,q)\) with
\(n\le X\) composite, \(q>Y\) prime and
\(q\mid\gcd(n,2^{n-1}-1)\) is
\[
                        \ll \frac{(X/Y)^2}{\log Y}.
\]
If \(n\ge5\) is odd, then \(n\) is prime if and only if
\(\gcd(n,2^{n-1}-1)\) has a prime factor exceeding
\(2n/\log_2 n\).
\end{theorem}

\begin{proof}
Write \(n=qk\). Since \(n-1=qk-1\equiv k-1\pmod{q-1}\),
Fermat's little theorem gives
\begin{equation}\label{eq:cofactor}
             q\mid 2^{n-1}-1\quad\Longleftrightarrow\quad
             q\mid 2^{k-1}-1.
\end{equation}
If \(n\) is composite, \(2\le k\le X/Y\). Distinct primes
\(q>Y\) dividing \(2^{k-1}-1\) have product at most \(2^{k-1}-1\).
Their number is at most \((k-1)\log 2/\log Y\). Summing over
\(k\) counts pairs and proves the first assertion.

If \(q>2n/\log_2 n\), then \(k<\tfrac12\log_2 n\).
Also \(2n/\log_2 n>\sqrt n\) for \(n\ge5\), as follows
from \(\log_2 n<2\sqrt n\).
For \(k\ge2\) this gives \(0<2^{k-1}-1<\sqrt n<q\), contrary
to \eqref{eq:cofactor}. Thus \(k=1\). Conversely, if \(n\ge5\)
is prime, \(q=n\) divides the gcd and exceeds the indicated
threshold.
\end{proof}

\section{A large prime divisor of the exponent}

The prime \(q\) in Theorem~\ref{thm:common} divides the modulus
\(n\). We now examine an independent condition on the exponent
\(n-1\).

\begin{lemma}\label{lem:exponent}
Let \(n\ge3\) and \(n-1=k\ell\), where \(\ell>\sqrt n\) is
prime. Then \(n\) is prime if and only if \(\ell\mid\varphi(n)\).
Equivalently, some residue \(u\not\equiv1\pmod n\) satisfies
\(u^\ell\equiv1\pmod n\). If \(n\) is composite, every integer
\(a\) satisfies
\[
    \gcd(n,a^{n-1}-1)=\gcd(n,a^k-1).
\]
\end{lemma}

\begin{proof}
If \(n\) is prime, its cyclic unit group has order \(k\ell\) and
contains an element of order \(\ell\). Conversely, any nontrivial
\(\ell\)-th root has order \(\ell\) in the group of units modulo
\(n\), so \(\ell\mid\varphi(n)\). Since \(\ell\nmid n\), the
product formula for \(\varphi(n)\) then gives a prime \(r\mid n\)
with \(r\equiv1\pmod\ell\). The cofactor \(n/r\) is also \(1\bmod
\ell\) and satisfies \(1\le n/r<\ell\). Thus \(n=r\).

Now let \(n\) be composite and let \(r^j\mid n\) be any prime
power. If \(r\mid a\), neither power of \(a\) in the displayed
gcd identity is congruent to \(1\bmod r\). Otherwise the order
of \(a\bmod r^j\) divides \(\varphi(r^j)\), which divides
\(\varphi(n)\). That order is coprime to \(\ell\), so it divides
\(k\ell\) exactly when it divides \(k\). This proves equality
of all prime-power divisibility conditions in the two gcds.
\end{proof}

\begin{theorem}\label{thm:joint}
Fix an integer \(a\ge2\), \(\theta>1/2\) and \(\delta>1-\theta\).
The number of composite \(n\in(X,2X]\) for which \(n-1\) has a
prime factor \(\ell>X^\theta\) and \(\gcd(n,a^{n-1}-1)\) has a
prime-power divisor at least \(X^\delta\) is
\[
  \ll_{a,\theta,\delta}
       \frac{X^{2-\theta-\delta}+X^{2-2\theta}}{\log X}.
\]
\end{theorem}

\begin{proof}
If \(\theta>1\), no such \(\ell\) exists for large \(X\). We may
therefore take \(1/2<\theta\le1\), so \(\delta>0\).
For sufficiently large \(X\), we have \(\ell>\sqrt{2X}\).
Write \(n=k\ell+1\), so \(k\le2X^{1-\theta}\). By
Lemma~\ref{lem:exponent}, the prime power in the hypothesis divides
\(a^k-1\). For each prime \(r\) for which some power at least
\(X^\delta\) divides \(a^k-1\), take the least such power \(s_r\).
The \(s_r\) are pairwise coprime, so their product is at most
\(a^k-1\). Their number is at most \(k\log a/(\delta\log X)\).

At least one of these \(s_r\) divides \(n\). For fixed \(k,s_r\),
the congruence \(k\ell\equiv-1\pmod{s_r}\) is impossible if
\(k\) and \(s_r\) have a common prime divisor. Otherwise it specifies
one residue class for \(\ell\). The condition \(X<k\ell+1\le2X\)
allows at most \(X/(k s_r)+1\le X^{1-\delta}/k+1\) integers
in that class. Summing this bound over the selected prime powers
and over \(k\le2X^{1-\theta}\) proves the theorem.
\end{proof}

\section{Consequences for prime pairs}

Averaging over the base gives a different measure of the contribution
from semiprimes. For odd \(n\ge3\), let
\[
   F(n)=\frac1n\#\{a\bmod n:a^{n-1}\equiv1\pmod n\}.
\]
Every residue satisfying the congruence is a unit. The denominator
counts all residue classes, including nonunits.

\begin{lemma}\label{lem:average}
The sum of \(F(n)\) over odd products of two primes \(n\le X\),
counting each integer once and allowing squares, is \(O(\sqrt X)\).
\end{lemma}

\begin{proof}
For \(n=rs\) with distinct odd primes, cyclicity of the unit groups
modulo \(r\) and \(s\) and the Chinese remainder theorem give
\(F(rs)=g^2/(rs)\), where \(g=\gcd(r-1,s-1)\).
Write \(r=ag+1\) and \(s=bg+1\) with \(a<b\). Then
\(F(rs)\le1/(ab)\) and \(g\le\sqrt{X/(ab)}\). Summing
over all positive \(a,b\) gives an upper bound of
\[
   \sqrt X\sum_{a,b\ge1}(ab)^{-3/2}\ll\sqrt X.
\]
Prime squares contribute at most \(\sqrt X\), since \(F(n)\le1\).
\end{proof}

Let \(T(X)\) count odd primes \(n\le X\) for which \(n-2\) is
prime. Let \(B(X)\) count odd integers \(n\le X\) having at most
two prime factors counted with multiplicity, with \(n-2\) prime
and \(n\mid 2^{n-1}-1\). Theorem~\ref{thm:semiprime} gives
\[
             0\le B(X)-T(X)\ll X^{3/4}(\log X)^{-5/4}.
\]
Likewise, let \(H_{\theta}(X)\) count odd \(n\le X\) with
\(n-2\) prime and a prime divisor \(q>X^\theta\) of
\(\gcd(n,2^{n-1}-1)\), where \(1/2<\theta<1\) is fixed.
For \(X^\theta\ge2\), Theorem~\ref{thm:common} gives
\[
 0\le H_{\theta}(X)-\bigl(T(X)-T(X^\theta)\bigr)
          \ll \frac{X^{2(1-\theta)}}{\theta\log X}.
\]
Every prime \(n>X^\theta\) counted by \(T(X)\) supplies the divisor
\(q=n\). Every other integer contributing to \(H_\theta(X)\)
is composite. Both upper bounds discard the condition that \(n-2\)
is prime when estimating this composite contribution.

For the corresponding average over bases, Lemma~\ref{lem:average}
and \(F(n)=1-1/n\) at an odd prime give
\[
  \sum_{\substack{3\le n\le X,\ n\ \mathrm{odd}\\
                n-2\ \mathrm{prime},\ \Omega(n)\le2}} F(n)
       =T(X)+O(\sqrt X),
\]
where \(\Omega(n)\) counts prime factors with multiplicity.
The discarded \(1/n\) terms contribute at most \(1+\log X\).

At the logarithmic cutoff of Theorem~\ref{thm:common}, one qualifying
integer \(n=p+2\) for arbitrarily large primes \(p\) would suffice
for infinitude. The divisor identity counts the exceptions when a
smaller cutoff is used. Proving the existence of qualifying divisors
at the logarithmic cutoff remains the missing step.


\setlength\emergencystretch{1.5em}
\section{A simultaneous order test}

Write \(L=\log X\) and \(z=\exp(\sqrt L)\). Define
\(E_n=(n^2-1)/2\) for odd \(n\).
\begin{theorem}\label{thm:simultaneous}
Among \(X<n\le2X\) with \(n\equiv5\bmod6\), the count satisfying
\(2^{E_n}\equiv1\bmod n(n+2)\) differs from the twin-prime count by
\(O(X\exp(-c\sqrt{\log X}))\), for an absolute \(c>0\).
If \(0<\beta<1/2\) is fixed and both sites have no prime factor
at most \(X^\beta\), the error is \(O_{\beta,\epsilon}(X^{1-\beta/3+\epsilon})\).
\end{theorem}
\begin{proof}
If both sites are prime then \(n-1,n+1\mid E_n\), so the congruence
holds. For a composite site \(pr\) with \(p\) prime and \(r>1\),
put \(k_p=\operatorname{ord}_p(2)\), so \(p\equiv1\bmod k_p\).
On the left \(k_p\mid r^2-1\) and on the right
\(k_p\mid(r-2)^2-1\). Either congruence has at most
\(4\cdot2^{\omega(k_p)}\) solutions modulo \(k_p\) by the
prime-power count for \(x^2=1\) and the Chinese remainder theorem.

For \(y\ge 3\) the number of \(p\ge y\) with \(k_p\le y^{1/3}\)
is \(O(y^{2/3}/\log y)\): all divide
\(\prod_{i\le y^{1/3}}(2^i-1)\). Their multiples up to \(3X\)
cost \(O(Xy^{-1/3}+y^{2/3})\) for \(y\le\sqrt{3X}\).
For the rough estimate use \(y=X^\beta\) and choose the smallest
prime factor of the composite site, at most \(\sqrt{3X}\).
For the remaining primes the count is
\[
 \ll_\epsilon\sum_{y<p\le\sqrt{3X}}k_p^\epsilon
               \left(\frac{X}{pk_p}+1\right)
 \ll_{\beta,\epsilon} X^{1-\beta/3+\epsilon},
\]
after harmlessly decreasing the exponent in the divisor bound.

For the unrestricted estimate, sites having only factors \(\le z\)
contribute \(O(Xe^{-c\sqrt L})\). Indeed Rankin's upper bound with
\(\sigma=1-1/(2\log z)\) is
\((3X)^\sigma\prod_{p\le z}(1-p^{-\sigma})^{-1}\); the product
is \(O((\log z)^C)\) by the prime harmonic estimate.
Choose any prime factor \(p>z\) and use the small-order estimate
with \(y=z\). For \(p\le X^{3/4}\) and \(k_p>z^{1/3}\), bound the
number of root classes by \(O(\sqrt{k_p})\) in the main term and by
\(O_\epsilon(X^\epsilon)\) in the rounding error. This costs
\(O(Xz^{-1/6}\log\log X+X^{3/4+\epsilon})\).
For \(p>X^{3/4}\) we have \(r\le3X^{1/4}\), and \(p\) divides
\(2^{|B(r)|}-1\) for \(B(r)=r^2-1\) or \((r-2)^2-1\).
The only positive roots are \(1\), and on the right also \(3\);
the right site is \(1\bmod6\), so \(r\ne3\). For each \(r\) the
nonzero integer has \(O(r^2+1)\) distinct prime factors. Summing
gives \(O(X^{3/4})\).
\end{proof}


\subsection{Several fixed bases}
Let \(b_1,\ldots,b_h\) be distinct fixed primes. Impose the congruence
\(b_i^{E_n}=1\bmod n(n+2)\) for all \(i\) on the same \(X^\beta\)-rough
rows as in Theorem~\ref{thm:simultaneous}. Apart from finitely many
small \(X\), all the bases are units. The number of false positives is
\begin{equation}\label{eq:multibase}
 O_{b_1,\ldots,b_h,\beta,\epsilon}
    (X^{1-\beta h/(h+1)+\epsilon}+X^{1/2+\epsilon}).
\end{equation}
To prove this bound, write \(K_p\) for the order of the subgroup
generated by the bases in \(\mathbb F_p^\times\). For
\(Y\le p<2Y\) the number of primes with \(K_p\le T\) is
\(O_b(T^{1+1/h}/\log Y)\).
Indeed, the monomials in the bases with exponent vectors from a cube
of side \(\lceil(2T)^{1/h}\rceil+1\) must collide modulo each such
prime. For each nonzero exponent difference cancel common factors;
the two resulting monomials are unequal integers by unique
factorization. There are \(O_b(T)\) differences, and the product of
their nonzero absolute differences has logarithm \(O_b(T^{1+1/h})\).
Decompose \(K_p\) dyadically up to \(Y^{h/(h+1)}\); for larger
orders use the ordinary prime-counting bound. This gives
\[
 \sum_{Y\le p<2Y}\frac1{K_p}\ll_b Y^{1/(h+1)}/\log Y.
\]
The common order divides \(E_n\) since \(\mathbb F_p^\times\) is cyclic.
For the smallest prime of a composite site, the same root-class
count as before, now modulo \(K_p\), is \(O_\epsilon(X^\epsilon)\).
Sum its main term \(X/(pK_p)\) over dyads \(p>X^\beta\), and its
unit rounding error over \(p\le\sqrt{3X}\), to obtain
\eqref{eq:multibase}.

\section{Two random bases and intersecting templates}
Assume \(1/3<\beta<1/2\). Continue to take \(n\equiv5\bmod6\) in the
dyad and retain \(P^-(n(n+2))>X^\beta\). Set \(M=n(n+2)\) and let
\[
 \kappa(n)=
 \frac{\#\{b\in(\mathbb Z/M)^\times:b^{E_n}=1\bmod M\}}{\varphi(M)} .
\]
Then \(\kappa(n)^2\) is the probability that two independent uniform
units both pass. Write \(T_X\) for the number of twin pairs starting
in the dyad. Every twin has \(\kappa(n)=1\). Each composite site is
a product of precisely two primes, counted with multiplicity, since
\(X^{3\beta}>2X+2\) eventually.
\begin{theorem}\label{thm:templates}
Remove each row on which either site has the form
\((at+1)(bt+1)\), for some positive integers \(a\le b\le X^{1/7}\)
and \(t\ge1\). On the remaining rough rows,
\[
 \sum_n\kappa(n)^2=T_X+O_{\beta,\epsilon}(X^{3/7+\epsilon}).
\]
Without this deletion, the weight of rough rows on which both sites
are composite is \(O_{\beta,\epsilon}(X^{3/7+\epsilon})\).
\end{theorem}
\begin{proof}
If \(p^2\mid M\), then \(p>3\) and \(p\nmid E_n\), since
\(n\equiv0\) or \(-2\bmod p\). Thus the cyclic local unit group
at \(p^2\) gives \(\kappa(n)\le1/p\).
At most \(O(\sqrt X)\) rows have a square site; their total weight
is \(O(X^{1/2-2\beta})\).

For a distinct-prime composite site \(pq\), set
\(p=ga+1,\ q=gb+1,\ g=(p-1,q-1)\), taking \(a\le b\) coprime.
Set \(u=(p-1)/(p-1,E_n)\), \(v=(q-1)/(q-1,E_n)\), and \(w=uv\).
A row with \(\kappa(n)\ge1/H\) has \(w\le H\), by the two local
cyclic group factors. With \(s=1\) for \(pq=n\), \(s=-1\) for
\(pq=n+2\), reduction of \(2E_n=n^2-1\) modulo \(p-1\) or \(q-1\)
gives
\[
    a\mid u(gb+2s), \qquad b\mid v(ga+2s).
\]
For example, modulo \(p-1\) the product remaining on the right of
\(n^2-1\) is \((q-1)(q+1)\) on the left site, and
\((q-1)(q-3)\) on the right site. After cancellation of \(g\)
one uses \((a,b)=1\). Multiplying the displayed congruences yields
the positive integer \(J=w(2g(a+b)+4s)/(ab)\). It satisfies
\[
          (Ja-2wg)(Jb-2wg)=4w(wg^2+sJ).
\]
The right side is nonzero: if \(s=-1,\ J=wg^2\), then
\((ga-2)(gb-2)=0\), whereas \(p,q>3\).

For fixed \(g,u,v\) there are \(O_\epsilon(X^\epsilon\sqrt{wg})\)
possibilities. When \(a\le\sqrt{wg}\), use \(b\mid v(ga+2s)\)
and the divisor bound; the right side of that divisibility is nonzero.
Otherwise \(1\le J\ll\sqrt{wg}\), and the factored identity gives
\(O_\epsilon(X^\epsilon)\) choices for each \(J\). All integers
whose divisors we count have polynomial size for \(H\le 2X^{2/7}\).
Summing over \(g\le G\) and \(uv\le H\) gives
\(O_\epsilon(X^\epsilon(HG)^{3/2})\).

For \(g>G\), necessarily \(ab\ll X/G^2\) and \(g\ll\sqrt{X/(ab)}\).
For fixed \(a,b,u,v\) the two congruences put \(g\) in at most one
class modulo \((a/(a,u))(b/(b,v))\). The total rounding error
is \(O_\epsilon(HX^{1+\epsilon}/G^2)\).
Expanding the gcds in divisor sums shows
\(\sum_{uv\le H}(a,u)(b,v)\ll_\epsilon H X^\epsilon\)
uniformly for \(a,b\ll X\). The total contribution from progression
lengths for \(b>Y\) is consequently
\[
 \ll_\epsilon X^\epsilon H\sqrt X
       \sum_{a\le b,\ b>Y}(ab)^{-3/2}
 \ll_\epsilon X^\epsilon H\sqrt{X/Y}.
\]
Choose \(G=X^{2/7}H^{-1/7}\), \(Y=X^{1/7}\).
Outside the deleted templates the number of false rows with
\(\kappa\ge1/H\) is \(O_\epsilon(X^{3/7+\epsilon}H^{9/7})\).
Sum in dyadic ranges of \(\kappa\ge X^{-2/7}\), and use the trivial
\(O(X\cdot X^{-4/7})\) below that cutoff. Twin pairs are never
deleted, which proves the first assertion.

For the second assertion keep templates with \(ab\le Z=X^{1/5}\)
on each composite site. The same proof outside these templates
uses \(\sum_{ab>Z}(ab)^{-3/2}\ll Z^{-1/2}\log(2Z)\) in place of
the sum with \(b>Y\). The length term is
\(O_\epsilon(HX^{2/5+\epsilon})\); the preceding squared-weight
bound still applies if either site is outside these templates.

For two retained templates write
\[
 (ch+1)(dh+1)-(at+1)(bt+1)=2,\quad A=ab,\quad B=cd .
\]
On putting \(U=2At+a+b,\ V=2Bh+c+d\), one obtains
\[
 BU^2-AV^2=D_0,\qquad
 D_0=B(a-b)^2-A(c-d)^2-8AB .
\]
We claim \(D_0\ne0\). Otherwise \(x=b/a,\ y=d/c\) satisfy
\(x+x^{-1}-y-y^{-1}=8\). Thus \(z=xy>0\) satisfies \(z\ne1\)
and \((x-y)(z-1)=8z\). It follows that
\[
 \big((x+y)(z-1)/2\big)^2=z(z^2+14z+1).
\]
With \(z=j-5\) this is \(Y^2=j^3-j^2-64j+220\), elliptic curve
24.a2. Its rational group has order four \cite{E24}. The four points
in the \(z\)-coordinate are \(\infty,(0,0),(1,4),(1,-4)\), giving the
contradiction.

For each fixed pair of templates there are \(O_\epsilon(X^\epsilon)\)
solutions. If \(AB\) is square, factor \((BU)^2-ABV^2=BD_0\) over
the integers. Otherwise \(BU+V\sqrt{AB}\) is an algebraic integer of
norm \(BD_0\ne0\) in a real quadratic field. Its ideal divides
\((BD_0)\), which has \(O_\epsilon(X^\epsilon)\) ideal divisors,
uniformly in the field by quadratic prime splitting. Generators of the
same ideal differ by units; both embeddings of a generator and their
reciprocals have sizes bounded by a fixed power of \(X\).
Every positive real quadratic unit greater than one is at least
\((1+\sqrt5)/2\), by its integral trace and norm. This leaves
\(O(\log X)\) possibilities for a generator, up to sign. Finally
there are \(O(Z^2\log^2(2Z))\) template pairs, giving
\(O_\epsilon(X^{2/5+\epsilon})\) rows in their intersections even
without the probability weight. This completes the proof.
\end{proof}

The second bound isolates the possible excess from a single
semiprime site beside a prime. The first bound requires a lower
estimate on its retained passing mass to give information on the
number of twin pairs.


\subsection{The untrimmed mixed rows}
\begin{corollary}\label{cor:mixed}
On all the rough rows of Theorem~\ref{thm:templates}, with no
template deletion,
\[
    0\le \sum_n\kappa(n)^2-T_X
       \ll_{\beta,\epsilon} \frac{\sqrt X}{(\log X)^{3-\epsilon}} .
\]
\end{corollary}
\begin{proof}
The composite/composite and square-site cases were already bounded.
Consider a remaining mixed row with composite site
\(m=(ag+1)(bg+1)\), \(a<b\) coprime and \(g\) the even gcd of the
two prime factors minus one. Put \(t=g/2\), and take \(s=1\) for a
left composite, \(s=-1\) for a right composite. The prime on the
other site supplies local passing probability one. Direct reduction
of \(E_n/g\) modulo \(a,b\) gives
\[
 \kappa(n)=\frac{(a,s+bt)(b,s+at)}{ab}.
\]
In fact \(m-1=g(gab+a+b)\), and the other factor in \(E_n/g\) is
\((m+1)/2\) or \((m-3)/2\), according to the orientation. Reduction
modulo \(a\) gives \(b(s+bt)\), with \(b\) a unit; the other
calculation is symmetric.

Write \(Z=X^{1/5}\). Outside \(ab\le Z\) the bound
\(O_\epsilon(X^{3/7+\epsilon})\) for one composite site from the
second part of Theorem~\ref{thm:templates} still applies.
For the retained templates sum over \(u,v\ge1\) with weights
\((uv)^{-2}\) and the necessary conditions
\[
 a\mid u(s+bt),\qquad b\mid v(s+at).
\]
This dominates the true \(\kappa(n)^2\): take \(u,v\) to be the two
local indices from the preceding proof. Each pair of conditions
allows at most one class of \(t\) modulo
\(Q=(a/(a,u))(b/(b,v))\). The interval has length
\(O(\sqrt{X/(ab)})\). Its rounding errors summed over \(ab\le Z\)
and \(u,v\) cost \(O(Z\log(2Z))\). Since
\[
 \sum_{u\ge1}(a,u)/u^2
 =\zeta(2)\sum_{h\mid a}\varphi(h)/h^2
 \le \zeta(2)\sigma_{-1}(a),
\]
the total length term with \(ab>Y=(\log X)^{12}\) is at most
\[
 \ll\sqrt X\sum_{ab>Y}(ab)^{-3/2}\sigma_{-1}(a)\sigma_{-1}(b)
 \ll\sqrt X\,Y^{-1/3}.
\]
The last estimate follows by inserting \((ab/Y)^{1/3}\) and
using the absolutely convergent Dirichlet series of \(\sigma_{-1}\)
at \(7/6\).

For \(ab\le Y\), consider each admissible class of \(t\).
Its parameter \(h\) satisfies \(t=t_0+Qh\), \(0\le t_0<Q\le ab\).
Count only parameters for which both linear polynomials \(2at+1\),
\(2bt+1\) and the quadratic
\[
       (2at+1)(2bt+1)+2s
\]
are prime. At the true sizes in our dyad \(t\asymp\sqrt{X/(ab)}\).
If the quadratic has square discriminant it factors in
\(\mathbb Z[t]\) up to an integer content. Its two nonconstant factors
may be taken with coefficients polynomial in \(a,b\) by the rational
root theorem. Both factors then have magnitude greater than one
at the true sizes for all large \(X\). A zero discriminant is also
excluded by \((a-b)^2=8s ab\), which has no positive rational ratio
\(a/b\). In the remaining case put \(D'=(a-b)^2-8s ab\) and let
\(\chi\) be the nonprincipal primitive quadratic character of its
squareclass. Outside the primes dividing \(2abQ(a-b)D'\), the number
of forbidden residues modulo \(p\) is \(3+\chi(p)\). Indeed the
two linear roots are distinct and neither annuls the quadratic.
The substituted discriminant has the same squareclass since \(Q\)
is invertible. Every excluded prime is bounded by a fixed power
of \(\log X\).

Fix \(0<\eta<1/10\) and a sufficiently small absolute \(\delta>0\).
Sieve by primes \(\exp(L^\eta)<p\le X^\delta\), with \(L=\log X\).
For these primes the root sum satisfies
\[
 \sum_p\frac{3+\chi(p)}p=3(1-\eta)\log L+O_\eta(1).
\]
Here prime harmonic summation treats the constant, while
Siegel--Walfisz, with partial summation, applies uniformly to the
nonprincipal character of conductor \(\ll L^{O(1)}\): throughout
this interval the conductor is \(\ll(\log p)^{O(1/\eta)}\).
The ordinary upper-bound sieve in dimension at most four
\cite[Chapter~19]{Kouk} gives
\[
 \ll_\eta \frac{\sqrt X}{\sqrt{ab}\,Q} L^{-3+3\eta}
             +O(X^{1/8})
\]
parameters. For the error take sieve level \(X^{1/10}\), choose
\(\delta\) small relative to that level, and use the CRT bound
\(O(4^{\omega(d)})\) for the remainder at each squarefree sieve
index \(d\). The standard multiplicative upper dimension bound
follows from the pointwise root bound four and Mertens' estimate.
The true prime factors all exceed \(X^\delta\) for large \(X\).
The power errors cost \(O(X^{1/8}Y\log(2Y))\) after summing the
convergent weights in \(u,v\). The length terms before the sieve
factor sum to \(O(\sqrt X)\) by the preceding divisor-series argument.
Take \(3\eta<\epsilon\) to conclude.
\end{proof}

There is also a high-order selection valid for every fixed \(\beta>0\)
in the rough range. For odd \(n\) put
\[
   \Delta=(n-1)(n+1)-\varphi(n)\varphi(n+2).
\]
This is zero exactly when both sites are prime. On any other rough
row it is positive and \(O_\beta(X^{2-\beta})\): each site has
\(O_\beta(1)\) factors, and the Euler product for its totient loses at
most \(O_\beta(X^{-\beta})\) proportionally. Every unit in the
passing subgroup has order dividing both \(2E_n\) and
\(\varphi(n(n+2))\), hence also \(\Delta\). Choosing
\(D=C_\beta X^{2-\beta}\) with sufficient \(C_\beta\) rules out
every composite row with a passing unit of order greater than \(D\).
At a twin the full unit group is \(C_2\times C_{E_n}\), since
\((n-1,n+1)=2\). At most \(2D\tau(E_n)\) elements have order at most
\(D\), by projection onto the cyclic factor. The fraction discarded
there is \(O_{\beta,\epsilon}(X^{-\beta+\epsilon})\).
The existence of arbitrarily large inputs carrying such a
high-order passing unit is the unresolved part.

\section{A shifted logarithmic gcd moment}

For \(n>1\) put \(G(n)=\log\gcd(n,2^{n-1}-1)\) and
\(\vartheta(n)=\mathbf1_{\{n\ \mathrm{prime}\}}\log n\).
Define a label \((p,u)\) by \(p\) odd prime and
\(1\le u\le v_p(2^{p-1}-1)\). Write \(A=p^uk_p\) and
\(w_A=\log p\), with \(k_p\) as in the preceding section.
For a second label \((q,v)\), put \(B=q^v k_q\).
Set
\[
 C_+=\sum_A\frac{w_A}{\varphi(A)},\qquad
 C_-=\sum_{3\nmid k_p}\frac{w_A}{\varphi(A)},\qquad
 C_0=\sum_{A,B}^{*}\frac{w_Aw_B}{\operatorname{lcm}(A,B)} .
\]
The star restricts the two labels by
\[
 p\ne q,\quad (p^u,k_q)=1,\quad(k_p,q^v)\mid3,\quad(k_p,k_q)\mid2.
\]
\begin{theorem}\label{thm:gcdmoment}
The three series converge absolutely. For every fixed \(J>0\),
\[
 \sum_{X<n\le2X}G(n)G(n+2)
 =
 \sum_{X<n\le2X}\vartheta(n)\vartheta(n+2)
 + X(C_0+C_++C_-)+O_J(X/(\log X)^J).
\]
\end{theorem}
\begin{proof}
For \(p\mid n\), \(p\ne2\), lifting the exponent gives
\(v_p(2^{n-1}-1)=v_p(2^{p-1}-1)\) if \(k_p\mid n-1\);
otherwise the valuation is zero. Here \(p\nmid n-1\) and
\(k_p\mid p-1\). Thus \(G(n)\) is the sum of \(w_A\) over the
conditions \(p^u\mid n,\ n\equiv1\bmod k_p\). Denote the part with
\(p^u\le z=\exp(\sqrt{\log X})\) by \(S_z(n)\).

Apart from \(O(Xe^{-c\sqrt L})\) exceptions in \(X<n\le2X+2\),
\(G(n)=S_z(n)+\vartheta(n)\). A composite with a contributing
\(p>z\) satisfies \(n=pr,\ k_p\mid r-1,\ r>1\).
The small-order product bound from Theorem~\ref{thm:simultaneous}
and the resulting single-class estimate for \(z<p\le X^{3/4}\)
give \(O(Xz^{-1/3}\log\log X+X^{3/4}+z^{2/3})\) such \(n\).
For \(p>X^{3/4}\), the integer \(2^{r-1}-1\ne0\) has \(O(r)\)
prime divisors and \(r\le3X^{1/4}\). Finally, sites with some
\(p^u>z,\ u\ge2,\ p^u\mid n\) cost
\[
 3X\sum_{p^u>z,\ u\ge2}p^{-u}
 \ll X z^{-1/4}\sum_{p,\,u\ge2}p^{-3u/4}\ll X z^{-1/4}.
\]
All the individual functions are \(O(L)\), which absorbs the
exceptional set into the stated error.

The constant tails are also small. For \(p\in[Y,2Y]\), the
small-order bound supplies \(O(Y^{2/3}/\log Y)\) primes with
\(k_p\le Y^{1/3}\); the rest have \(1/k_p\le Y^{-1/3}\).
Using \(A/\varphi(A)\ll\log\log(3A)\), the first-power sum
\(\sum_{Y<p\le2Y}(\log p)/\varphi(pk_p)\) is
\(O(Y^{-1/3}\log^2(2Y))\).
For higher powers \(\sum_{p^u>z,\,u\ge2}w_A/\varphi(A)
\ll z^{-1/8}\), by the same convergent prime-power bound with
logarithms absorbed (use \(k_p\le p\)).
In the starred sum \((A,B)\le6\) by its displayed conditions.
These estimates prove convergence and permit replacing each
truncated sum by its full value with all-logarithmic precision.

Expand the product of the two short sums. The simultaneous
congruences on \(n\) are \(0,1,-2,-1\) modulo \(p^u,k_p,q^v,k_q\),
respectively. Their compatibility is exactly the starred condition.
The Chinese remainder theorem then gives
\(X/\operatorname{lcm}(A,B)+O(1)\) occurrences.
There are at most \(z^2\) ordered pairs of retained labels, and their
rounding errors are negligible. For
\(\vartheta(n)S_z(n+2)\) the prime is in the class \(-2\bmod p^u\),
\(-1\bmod k_p\), always primitive. For
\(S_z(n)\vartheta(n+2)\) the prime \(n+2\) has residues \(2,3\)
respectively, giving a primitive class exactly if \(3\nmid k_p\).
All moduli \(A\le z^2\) are distinct: the largest prime factor
of \(A\) is \(p\), with exponent \(u\). Bombieri--Vinogradov
\cite{BV} for the log-weighted prime sums, with partial summation
or the maximal-interval form at the shifted endpoints, yields
\(X C_+\) and \(X C_-\) up to the stated error. Replacing the
von Mangoldt weights in this theorem by primes costs at most
\(O(X^{1/2+o(1)}z^2)\) on the short list of moduli.
\end{proof}

The three constants isolate the composite and mixed terms in this
logarithmic statistic. A lower bound exceeding their contribution
would still have to come from the correlation at the actual shift.

\section*{Acknowledgments}
ChatGPT and Codex contributed to the mathematical derivations,
drafting and proof checks.

\begin{thebibliography}{9}
\bibitem{LRJ}
F.~Luca, D.~Ralaivaosaona and J.~Jim\'enez Urroz,
\emph{On pseudoprime RSA moduli}, Afrika Matematika
\textbf{36} (2025), article 118.
\href{https://doi.org/10.1007/s13370-025-01338-1}{doi:10.1007/s13370-025-01338-1}.

\bibitem{MV}
H.~L.~Montgomery and R.~C.~Vaughan,
\emph{The large sieve}, Mathematika \textbf{20} (1973), 119--134.
\href{https://doi.org/10.1112/S0025579300004708}{doi:10.1112/S0025579300004708}.
\bibitem{BV} E. Bombieri, \emph{On the large sieve},
Mathematika \textbf{12} (1965), 201--225.
\href{https://doi.org/10.1112/S0025579300005313}{Author's paper}.

\bibitem{E24} The LMFDB collaboration, \emph{Elliptic curve 24.a2
(Cremona label 24a3)}, \href{https://www.lmfdb.org/EllipticCurve/Q/24/a/2}{curve equation and Mordell--Weil group}.
\bibitem{Kouk} D. Koukoulopoulos, \emph{The Distribution of Prime
Numbers}, Graduate Studies in Mathematics \textbf{203}, American
Mathematical Society, 2019. \href{https://dms.umontreal.ca/~koukoulo/documents/publications/primes.pdf}{Author's text}.
\end{thebibliography}
\end{document}
